\documentclass[../numerical_methods.tex]{subfiles}
\begin{document}
\section{Aula 10 - 09 de Abril, 2026}
\subsection{Motivações}
\begin{itemize}
	\item Computação Científica;
	\item Erros de Arredondamento.
\end{itemize}
\subsection{Computação Científica}
No computador, a memória é finita, então não é possível representar um número, como uma dízima não periódica, por completo. O padrão IEEE 754 (\textit{Double Precision}) usa 64 bits para representar um número real: 1 bit para o sinal, 11 para o expoente e 52 para a \textit{mantissa}, que é basicamente os dígitos significativos.
Além disso, frações simples na base 10 podem virar dízimas periódicas na base 2, forçando o computador a "cortar" o número, ou seja, o computador não entende, por exemplo, "0,1"; ao invés disso, ele entende bits, e o padrão IEEE 754 é a norma mundial para aritmética de ponto flutuante.

Imagine um número real x representado como:
\[
	x = \pm M \times B^{E}.
\]
Vamos explicar o que é cada parte:

\textbf{O Sinal}: 1 bit (0 para positivo, 1 para negativo).

\textbf{Mantissa} (M): É onde reside a precisão. No formato de precisão dupla (double), temos 52 bits. Se um número precisa de 53 bits para ser representado com exatidão, o 53º é "jogado fora", gerando o erro.

\textbf{Expoente} (E): Define a ordem de grandeza. É o que permite o computador representar desde o tamanho de um átomo até a distância entre galáxias usando o mesmo número de bits.

\begin{tcolorbox}[
		skin=enhanced,
		title=Observação,
		fonttitle=\bfseries,
		colframe=black,
		colbacktitle=cyan!75!white,
		colback=cyan!15,
		colbacklower=black,
		coltitle=black,
		drop fuzzy shadow,
		%drop large lifted shadow
	]
	Como a base é B=2, números que são frações exatas em base 10 (como \(\frac{1}{10}=0.1\)) tornam-se dízimas periódicas em binário (0.00011001100...2); como a mantissa é finita, o corte é inevitável.
\end{tcolorbox}
\begin{listing}[H]
	\begin{minted}[bgcolor = DarkBackground, linenos=true, breaklines]{Python}
# DEMONSTRAÇÃO 1: A ilusão da matemática de ponto flutuante
print("--- Teste de Soma Simples ---")

resultado = 0.1 + 0.2
print(f"O computador diz que 0.1 + 0.2 = {resultado}")

if resultado == 0.3:
    print("Conclusão: A matemática funciona perfeitamente!")
else:
    print("Conclusão: Ocorreu um erro de representação (arredondamento)!")

# Mostrando o que o computador realmente guardou na memória (com 55 casas decimais)
print("\n--- O que está escondido na memória ---")
print(f"Valor real de 0.1: {0.1:.55f}")
print(f"Valor real de 0.2: {0.2:.55f}")
print(f"Valor real de 0.3: {0.3:.55f}")
\end{minted}
\end{listing}

Quando o número real possui mais algarismos significativos do que a memória permite guardar, acontece um erro de armazenamento, conhecido também como \textbf{erro de arredondamento}; um exemplo disso é o erro de arrendondamento ao tentar armazenar \(\pi \) com apenas 4 casas -- ao invés de \(\pi = 3.141592\), fica \(\pi \approx 3.1416\).

Outro possível erro é quando substituímos um processo matemático finito, como uma aproximação por série de Taylor, por uma versão finita sua, resultando no erro metodológico conhecido como \textbf{erro de truncamento}, onde ``truncar'' é um verbo que dá a ideia de ``cortar''. No exemplo das séries de Taylor mencionadas, se considerarmos a de \(e^{x}\), ela seria
\[
	e^{x} = 1 + x + \frac{x^{2}}{2!} + \dotsc ,
\]
mas se somarmos apenas os 3 primeiros termos, o erro entre o resultado real e o calculado é o erro de truncamento.

Com base nos erros mencionados, dá para imaginar que o conjunto dos números da máquina não é um contínuo, tampouco estão igualmente espaçados -- na reta real, entre quaisquer dois números, sempre há um terceiro, mas isso não é verdade no computador! O conjunto de pontos flutuantes, \(\mathbb{F}\), é finito, e quebrar isso levanta ainda outros tipos de erros:
\begin{itemize}
	\item \textbf{Overflow}: ocorre quando o expoente E ultrapassa o valor máximo permitido pela arquitetura do computador: se ela libera 11 bits para precisão dupla, são permitidos expoentes até 1023, e se uma operação resulta em \(2^{1024}\), o computador não tem como registrar a magnitude -- ele desiste e retorna Inf (infinito);
	\item \textbf{Underflow}: ocorre quando o número é tão, tão pequeno (próximo de zero) que o expoente atinge seu limite inferior: se ele for o expoente -1022, o expoente -1023 seria arredondado para 0.
\end{itemize}
Com isso, faz sentido definir que
\begin{def*}
	O \textbf{épsilon \(\mathbf{\varepsilon }\)} de uma máquina é o limite superior do erro \textit{relativo} da representação, a distância entre o número 1.0 e o próximo número imediatamente representável na máquina. \(\square\)
\end{def*}
Para pensar nele, ele é o número tal que, se somarmos 1.0 a qualquer valor menor que, por exemplo, \(\varepsilon/2\), então o computador ``absorve'' o valor e ele continua sendo 1.0.

\begin{listing}[H]
	\begin{minted}[bgcolor = DarkBackground, linenos=true, breaklines]{Python}
import sys

# DEMONSTRAÇÃO 2: Encontrando o Épsilon da Máquina ao vivo
eps = 1.0

# O laço roda enquanto o computador conseguir notar a diferença entre 1 e (1 + eps/2)
while 1.0 + (eps / 2.0) > 1.0:
    eps = eps / 2.0

print(f"Épsilon calculado manualmente: {eps}")
print(f"Épsilon oficial do Python:   {sys.float_info.epsilon}")

# Provando o limite da máquina
print("\nO que acontece se somarmos metade do Épsilon ao número 1?")
numero_minusculo = eps / 2.0

if 1.0 + numero_minusculo == 1.0:
    print("O computador ignorou o número! 1.0 + minúsculo = 1.0 (Absorção)")
else:
    print("O computador conseguiu registrar a soma.")
\end{minted}
\end{listing}

No computador, o padrão de 64 bits consegue armazenar apenas em torno de 15 a 17 dígitos significativos, e isso dá origem a um fenômeno conhecido como \textbf{Cancelamento Catastrófico}, que é quando tentarmos subtrair dois números de grandezas muito próximas, como calcular \(\sqrt[]{x+1} - \sqrt[]{x}\) para \(x=10^{16}\); neste caso, os primeiro 15 dígitos das duas raízes são idênticos, e, ao subtraí-los, eles basicamente se anular! Consequentemente, a informação matemática é destruída e o computador preenche o vazio com zeros ou ``lixos de memórias''.

Para evitar essa perda de dados, é comum racionalizar a equação multiplicando pelo seu conjugado, que transforma a subtração problemática no numerador em uma adição segura no \textit{denominador}:
\[
	y = \sqrt[]{x+1}-\sqrt[]{x} \rightsquigarrow \overline{y} = (\sqrt[]{x+1}-\sqrt[]{x}) \frac{\sqrt[]{x+1}+\sqrt[]{x}}{\sqrt[]{x+1}+\sqrt[]{x}} = \frac{1}{\sqrt[]{x+1}+\sqrt[]{x}}
\]
A importância prática disso se dá no fato de que, se essa equação instável for o denominador de um algoritmo maior, por exemplo uma força \(F = \frac{100}{y}\), então o ``quase zero'' vira um zero absoluto, causando um erro fatal de divisão por zero e travando o software.

\begin{listing}[H]
	\begin{minted}[bgcolor = DarkBackground, linenos=true, breaklines]{Python}
import math

x = 10**16 # valor extremo

print(f"--- SIMULAÇÃO: Calculando F = 100 / y para x = {x} ---\n")

# --- CENÁRIO 1: Abordagem Direta (Fórmula Instável) ---
print("MÉTODO 1: Fórmula Direta (com subtração)")
try:
    # A máquina tenta subtrair números idênticos. O resultado vira 0.0
    y_instavel = math.sqrt(x + 1) - math.sqrt(x)

    # O programa tenta dividir por zero e quebra
    F_instavel = 100 / y_instavel
    print(f"Resultado F: {F_instavel}")

except ZeroDivisionError:
    print("ERRO FATAL: Divisão por zero!")
    print("   O y instável colapsou para 0.0 devido ao cancelamento catastrófico e quebrou o sistema.\n")

# --- CENÁRIO 2: Abordagem Algébrica (Fórmula Estável) ---
print("MÉTODO 2: Fórmula Racionalizada (com soma)")
try:
    # Trocamos a subtração por uma adição. A informação é preservada (5e-09)
    y_estavel = 1 / (math.sqrt(x + 1) + math.sqrt(x))

    # A divisão ocorre perfeitamente
    F_estavel = 100 / y_estavel
    print(f" SUCESSO: Força F = {F_estavel:.2f}")
    print("   A precisão foi mantida. O algoritmo sobreviveu e calculou corretamente.")

except Exception as e:
    print(f"Erro: {e}")
\end{minted}
\end{listing}

\subsection{Exercícios Teóricos e Práticos}

\textbf{A Origem dos Erros}: Diferencie claramente o Erro de Arredondamento do Erro de Truncamento em métodos numéricos. Dê um exemplo clássico para cada um dos casos.

\textbf{O Paradoxo da Soma}: Ao calcular a soma harmônica
\[
	S = \sum\limits_{n=1}^{N}\frac{1}{n}
\]
para um N muito grande (ex: 10 milhões) em um computador com precisão simples, um programador nota que calcular a soma de n=1 até N (ordem crescente de n) resulta em um valor ligeiramente diferente de calcular de n=N até 1 (ordem decrescente). Explique por que essa diferença ocorre e qual dos dois métodos produz o resultado mais preciso.

\textbf{Fugindo do Cancelamento Catastrófico} Considere a função matemática abaixo para valores de x muito próximos de zero (ex: \(x=10^{-8}\)):
\[
	f(x) = \frac{1-\cos^{}{(x)}}{x^{2}}.
\]

a) Explique por que calcular essa função diretamente no computador causará uma perda severa de algarismos significativos.
b) Utilizando a identidade trigonométrica \(\displaystyle \sin^{2}{(x)} = \frac{1-\cos^{}{(2x)}}{2}\), proponha uma equação matematicamente equivalente para f(x) que seja numericamente estável para valores pequenos de x.

\textbf{Racionalização de Raízes}: A função \(g(x) = \sqrt[]{x^{2}-1}-x\) sofre de instabilidade numérica quando avaliada para valores de x muito grandes. Reescreva a função g(x) de modo a evitar o cancelamento catastrófico.

\textbf{O Limite de Taylor e o Épsilon}: Sabemos que o número neperiano e pode ser calculado pela Série de Taylor infinita:
\[
	e = \sum\limits_{n=0}^{\infty}\frac{1}{n!} = 1 + 1 + \frac{1}{2!} + \frac{1}{3!}+\frac{1}{4!} + \dotsc
\]
Escreva um programa em Python que calcule o valor de e iterativamente. O programa deve parar de somar novos termos quando o próximo termo a ser adicionado, \(\frac{1}{n!}\), for tão pequeno que não fará mais diferença na soma total (ou seja, quando o termo for menor ou igual ao Épsilon da máquina, ou quando o valor da soma parar de mudar).

\subsubsection{Gabarito}
Respostas Teóricas e Conceituais

1. Resposta:

Erro de Arredondamento: É causado pela limitação física do hardware em representar números reais com infinitas casas decimais usando uma quantidade finita de memória (bits).

Exemplo: Armazenar o número \(\pi \) como 3.14159265 em uma variável ou tentar representar a fração 1/10 em base binária, o que gera uma dízima periódica que precisa ser cortada.

Erro de Truncamento: É originado pelo método matemático escolhido. Ocorre quando substituímos um procedimento infinito por um finito, independentemente do computador ser perfeito ou não.

Exemplo: Calcular o valor de sin(x) usando apenas os 3 primeiros termos de sua Série de Taylor, ignorando o restante da série infinita.

2. Resposta: Ocorre o fenômeno do acúmulo de erro por absorção. A soma mais precisa é a decrescente (de n=N até 1). Quando n é muito grande, o termo 1n é um número muito pequeno. Se somarmos do maior para o menor (começando por 11+12…), a variável acumuladora rapidamente se torna um número consideravelmente grande. Quando o laço chegar nos valores muito pequenos de \(\frac{1}{n}\), a diferença de magnitude entre a soma acumulada e o novo termo será tão grande que a mantissa não terá espaço para acomodar o termo pequeno, ocorrendo o underflow relativo (o computador ``ignora'' a soma e S+pequeno=S). Somando de n=N até 1, começamos somando números minúsculos entre si, permitindo que eles se acumulem até formarem um valor significativo antes de serem somados aos termos maiores.
Respostas Algébricas

3. Resposta: a) Para \(x\approx 0\), o valor de cos(x) se aproxima muito de 1. A operação \(1-\cos^{}{(x)}\) no numerador resultará na subtração de dois números quase idênticos (ex: 1 - 0.99999999), causando cancelamento catastrófico. Os dígitos significativos se anulam, restando apenas o erro de arredondamento.

b) Se \(\displaystyle \sin^{2}{(x)}=\frac{1-\cos^{}{(2x)}}{2}\), então podemos fazer uma substituição de variável (\(u = 2x\) é o mesmo que \(x = \frac{u}{2}\)):
\[
	1-\cos^{}{u} = 2\sin^{2}{}\biggl(\frac{u}{2}\biggr)
\]
Substituindo na equação original:
\[
	f(x) = \frac{2\sin^{2}{}\biggl(\frac{x}{2}\biggr)}{x^{2}}.
\]
Como não há subtração de termos próximos nesta nova fórmula, ela é numericamente estável para \(x\approx 0\).

4. Resposta: Para x muito grande, \(x^{2}+1 \approx x^{2}\), logo \(\sqrt[]{x^{2}+1}\approx x\), gerando a subtração perigosa x-x. Para resolver, multiplicamos a expressão pelo seu conjugado:
\begin{align*}
	 & g(x) = (\sqrt[]{x^{2}+1} - x) \frac{\sqrt[]{x^{2}+1}+x}{\sqrt[]{x^{2}+1}+x} \\
	 & g(x) = \frac{(x^{2}+1)-x^{2}}{\sqrt[]{x^{2}+1}+x}                           \\
	 & g(x) = \frac{1}{\sqrt[]{x^{2}+1}+x}
\end{align*}
Agora temos apenas uma soma no denominador, eliminando o risco de cancelamento catastrófico
\subsection{Calculando a Variância}
O cancelamento catastrófico afeta também a análise de dados, e pode destruí-las completamente se não for lidado corretamente.

Nesse sentido, a variância \(\sigma^{2}\) mede o quão distantes da média os dados estão, e tem uma de suas formas dada por
\[
	\mathrm{Var}(X) = \frac{\sum\limits_{i=0}^{N}x_{i}^{2}}{N} - \mu^{2}.
\]

Imagine que você está analisando dados com valores muito altos, mas com diferenças minúsculas, como 100000000001, 100000000002, 100000000003; a média ao quadrado e a soma dos quadrados serão dois números colossais e quase idênticos, tal que, ao subtrair um do outro, a precisão desaparece! O resultado pode ser nulo ou até mesmo negativo!

Por isso, bibliotecas como o Pandas e o NumPy nunca usam a fórmula clássica dos livros, usando algoritmos estáveis para proteger os dados, como o \textit{Algoritmo de Welford}.

\begin{listing}[H]
	\caption{Calculando a Variância}
	\begin{minted}[bgcolor = DarkBackground, linenos=true, breaklines]{Python}
import numpy as np

# --- O Conjunto de Dados ---
# Números muito grandes com diferenças muito pequenas
dados = [10**9 + 1.0, 10**9 + 2.0, 10**9 + 3.0]
N = len(dados)

print("Dados analisados:", dados)
print("-" * 50)

# --- MÉTODO 1: A "Fórmula do Livro" (Instável) ---
# Var(X) = E[X^2] - (E[X])^2
soma_quadrados = sum(x**2 for x in dados)
media = sum(dados) / N

variancia_instavel = (soma_quadrados / N) - (media**2)

print("MÉTODO 1 (Fórmula Clássica Escolar):")
print(f"Resultado da Variância: {variancia_instavel}")
if variancia_instavel < 0:
    print(" ERRO FATAL: Variância DEU NEGATIVA! Matematicamente impossível.")


print("\n" + "-" * 50 + "\n")


# --- MÉTODO 2: O Jeito Profissional (Estável - NumPy) ---
# O NumPy utiliza algoritmos estáveis nos bastidores
variancia_estavel = np.var(dados)

print("MÉTODO 2 (Algoritmo Estável do NumPy):")
print(f"Resultado da Variância: {variancia_estavel}")
print(f" SUCESSO: A variância real é exatamente {variancia_estavel}")
\end{minted}
\end{listing}
\end{document}
